\documentclass[twocolumn, 10pt]{article}
\usepackage[utf8]{inputenc}
\usepackage{amsmath}
\usepackage{authblk}
\usepackage{geometry}
\geometry{margin=0.75in}

\title{\textbf{Architectural Synthesis of Hybrid Instruction Set Computing (HISC): Leveraging Heterogeneous Pipelines for Energy-Efficient Edge Computing}}
\author{Castor \textit{et al.}}
\affil{Advanced Computing Research Initiative}
\date{March 2026}

\begin{document}

\maketitle

\begin{abstract}
This paper introduces and analyzes the Hybrid Instruction Set Computing (HISC) architecture, a novel processor design that dynamically routes instructions between RISC and CISC processing units[cite: 2]. By utilizing a specialized split L1 cache hierarchy and a robust Hardware Abstraction Layer (HAL), HISC optimizes for both high-frequency atomic operations and complex macro-instructions[cite: 3, 185]. We explore the transition of HISC from high-performance computing to power-constrained environments such as mobile and IoT, demonstrating its potential to coexist with future quantum paradigms[cite: 140, 163].
\end{abstract}

\section{Introduction}
The evolution of processor design has historically vacillated between the streamlined efficiency of Reduced Instruction Set Computing (RISC) and the semantic density of Complex Instruction Set Computing (CISC)[cite: 1, 189]. Modern computational demands, particularly in AI and real-time data analysis, require a synthesis of these philosophies[cite: 6]. HISC proposes a unified silicon footprint where instructions are intelligently routed to the optimal pipeline based on computational complexity[cite: 2, 17].

\section{Literature Review}
The philosophical precursors to HISC include ARM's \textit{big.LITTLE} architecture, which employs power-efficient and high-performance cores [cite: 15], and heterogeneous computing models that pair CPUs with GPUs[cite: 16]. However, HISC diverges by integrating disparate instruction set architectures (ISAs) into a single, dynamically scheduled execution environment[cite: 17, 31].

\section{Hypotheses}
\begin{itemize}
    \item \textbf{H1}: A split L1 cache architecture reduces inter-pipeline interference and allows for ISA-specific latency tuning[cite: 79, 120].
    \item \textbf{H2}: A Hardware Abstraction Layer (HAL) effectively masks ISA complexity, ensuring software portability across the hybrid core[cite: 170, 186].
    \item \textbf{H3}: HISC provides superior power-per-watt metrics in mobile SoCs compared to monolithic ISA designs[cite: 142].
\end{itemize}

\section{Methodology & Architectural Design}

\subsection{Microarchitecture and Pipeline Routing}
The HISC core employs an \textit{Instruction Profiler} to monitor sequences at the firmware level[cite: 32]. Simple arithmetic operations are shuttled to a short RISC pipeline, while complex string manipulations or specialized DSP tasks are routed to a microcoded CISC pipeline[cite: 35, 109].

\subsection{Memory and Cache Hierarchy}
To prevent "toe-stepping" between pipelines, we implement:
\begin{itemize}
    \item \textbf{Split L1 Caches}: The RISC L1 utilizes 32-byte lines for low-latency, high-frequency fetches, while the CISC L1 utilizes 64-byte lines for instruction density[cite: 89, 91].
    \item \textbf{Unified L2/L3}: A centralized coherence engine manages shared data, facilitating seamless collaboration between pipelines without L1 thrashing[cite: 82, 114].
\end{itemize}

\subsection{Software Stack (HAL)}
The HAL serves as the critical intermediary, defining generic APIs for system operations[cite: 177]. This allows the OS to see a uniform execution target while the hardware manages the underlying dispatch logic[cite: 49, 186].

\section{Results and Outcomes}
Initial theoretical modeling suggests that:
\begin{itemize}
    \item \textbf{Efficiency}: The split L1 design provides optimized latency for distinct fetch patterns, maintaining high IPC for RISC workloads[cite: 108, 119].
    \item \textbf{Flexibility}: Compilers (e.g., LLVM-based) can generate "fat" binaries containing routines optimized for both pipelines, selected at runtime based on power state[cite: 44, 46].
    \item \textbf{Market Viability}: In the IoT and mobile sectors, HISC offers a unique niche by providing a "brawny" pipeline for bursty encryption or AI tasks alongside a "lean" pipeline for background polling[cite: 150, 156].
\end{itemize}

\section{Conclusion & Further Work}
HISC represents a significant leap in classical architecture, particularly relevant for edge devices where quantum computing remains impractical[cite: 137, 164]. While manufacturing complexity remains a challenge, the potential for power-efficient, specialized computing is immense[cite: 7, 142].

\textbf{Further Work}:
\begin{enumerate}
    \item \textbf{Software Simulation}: Extending QEMU or gem5 to model dual-pipeline dispatch[cite: 56].
    \item \textbf{FPGA Prototyping}: Validating the HAL and split-cache coherence in Verilog/VHDL[cite: 58, 59].
    \item \textbf{Ecosystem Development}: Creating a Linux-based kernel with HISC-aware drivers[cite: 63, 64].
\end{enumerate}

\section*{Final Thoughts}
The transition from theoretical HISC to physical silicon requires an iterative approach—starting with simulation and moving toward ASIC tape-out[cite: 55, 60]. By focusing on the mobile and embedded markets, HISC can carve out a long-term coexistence strategy with emerging high-level computing paradigms[cite: 163, 165].

\end{document}